\subsection{Logic Evaluation Architecture}

Branching syntax defines where the possible routes are. Logic evaluation decides which route is selected. In Python, the expression written after \texttt{if}, \texttt{elif}, or \texttt{while} does not always need to be a literal \texttt{True} or \texttt{False}. Python can ask many kinds of objects whether they should count as true or false in a control-flow decision.

This is why Python conditionals are not limited to direct comparisons such as \texttt{x > 0}. They may contain names, containers, strings, numbers, function results, logical operators, and chained comparisons.

\begin{quote}
A conditional expression does not merely calculate a value. In a branch or loop header, Python converts that value into a routing decision.
\end{quote}

This subsection studies three parts of that routing machinery: truth-value testing, short-circuit evaluation, and comparison chaining.

\subsubsection{Truth Value Testing: Evaluating Implicit Truthiness and Falsiness via \texttt{\_\_bool\_\_} and \texttt{\_\_len\_\_}}

Python accepts any object in a conditional position. It then asks whether that object should be treated as true or false. This process is called truth-value testing.

The direct conversion can be seen with \texttt{bool()}:

\begin{verbatim}
values = [
    True,
    False,
    42,
    0,
    "spam",
    "",
    [1, 2, 3],
    [],
    None,
]

for value in values:
    print(f"{repr(value):>10} -> {bool(value)}")
\end{verbatim}

Possible output:

\begin{verbatim}
      True -> True
     False -> False
        42 -> True
         0 -> False
    'spam' -> True
        '' -> False
 [1, 2, 3] -> True
        [] -> False
      None -> False
\end{verbatim}

This does not mean that all these objects are booleans. It means they can be tested as booleans.

%<<<PROGRAM:programs/log07>>>
A list does not normally define its own \texttt{\_\_bool\_\_()} method. Instead, Python can use its length. An empty list is false. A non-empty list is true.

%<<<PROGRAM:programs/log06>>>
The practical rule is:

\begin{quote}
If an object has \texttt{\_\_bool\_\_()}, Python uses it for truth-value testing. If it does not, but it has \texttt{\_\_len\_\_()}, zero length means false and non-zero length means true.
\end{quote}

Numbers usually have direct boolean behavior:

%<<<PROGRAM:programs/log05>>>
Strings are length-based:

%<<<PROGRAM:programs/log04>>>
An empty string is false because its length is zero:

\begin{verbatim}
quote = ""

if quote:
    print("quote exists")
else:
    print("quote is empty")
\end{verbatim}

Possible output:

\begin{verbatim}
quote is empty
\end{verbatim}

This is common Python style. Instead of writing:

\begin{verbatim}
if len(quote) > 0:
    print("quote exists")
\end{verbatim}

Python code often writes:

\begin{verbatim}
if quote:
    print("quote exists")
\end{verbatim}

The shorter version does not mean that \texttt{quote} is a boolean. It means the string object is being used in a truth-value test.

The object \texttt{None} is always false:

%<<<PROGRAM:programs/log03>>>
The \texttt{is None} test checks identity with the single \texttt{None} object. It is not the same as a general false test.

%<<<PROGRAM:programs/log02>>>
The integer \texttt{0} is false in a truth-value test, but it is not \texttt{None}. This distinction matters in real programs. A missing value and a present-but-zero value are not the same fact.

\subsubsection{Short-Circuit Evaluation Mechanics and Object Return Behaviors}

Unlike C's \texttt{\&\&} and \texttt{||}, which normalize to integer status flags, Python's \texttt{and} and \texttt{or} are expression operators that return a \texttt{PyObject*} reference---the last operand actually evaluated---not a dedicated boolean register value. In a branch header the returned object is truth-tested afterward; in assignment the object identity and type of the result depend on which side short-circuited.

The compiler does not emit a separate boolean temporary for every logical chain. It fuses control flow into conditional jumps on the evaluation stack:

\begin{itemize}
    \item \texttt{JUMP\_IF\_FALSE\_OR\_POP}: if the TOS object is false, jump to the short-circuit target \emph{without} popping; if true, pop it and continue evaluating the right operand.
    \item \texttt{JUMP\_IF\_TRUE\_OR\_POP}: symmetric early exit for \texttt{or} when the left operand is already true.
\end{itemize}

Conceptual bytecode for \texttt{a and b}:

\begin{verbatim}
LOAD_NAME            a
JUMP_IF_FALSE_OR_POP exit_and    # false -> exit with a still on stack
LOAD_NAME            b
exit_and:
\end{verbatim}

For \texttt{a or b}, the jump polarity reverses: a true left operand skips evaluation of \texttt{b} and leaves that object as the expression result.

The operand-return semantics matter at the API boundary:

\begin{verbatim}
result = "" or "fallback"      # str 'fallback', not bool True
result = "ready" and "go"       # str 'go'
result = 0 and 42               # int 0 (short-circuited before 42)
\end{verbatim}

Guard idioms such as \texttt{items and items[0] == target} rely on the jump: an empty list is falsy, the \texttt{JUMP\_IF\_FALSE\_OR\_POP} fires, and the index operation never executes---not because of a special-case type rule, but because the right-hand subgraph is unreachable in the control-flow graph.

\begin{quote}
Short-circuit logical operators are stack-directed conditional jumps that leave an operand object on the stack, not boolean normalizers.
\end{quote}

\subsubsection{Comparison Chaining: Interpreting Expressions Such as \texttt{a < b < c}}

Python allows comparison expressions to be chained. The expression:

\begin{verbatim}
a < b < c
\end{verbatim}

means:

\begin{verbatim}
a < b and b < c
\end{verbatim}

but with one important structural detail: the middle expression is evaluated only once.

A simple example:

%<<<PROGRAM:programs/log01>>>
The result of the whole comparison chain is a \texttt{bool} object. The comparison itself still follows short-circuit logic. If the first comparison fails, Python does not need the later comparisons.

\begin{verbatim}
a = 30
b = 20
c = 10

check = a < b < c

print(f"check: {check}")
\end{verbatim}

Possible output:

\begin{verbatim}
check: False
\end{verbatim}

The first comparison \texttt{a < b} is false, so the whole chain is false.

This syntax is especially useful for interval tests:

\begin{verbatim}
value = 42

if 0 <= value <= 100:
    print("value is inside the accepted interval")
else:
    print("value is outside the accepted interval")
\end{verbatim}

Possible output:

\begin{verbatim}
value is inside the accepted interval
\end{verbatim}

Written without chaining, the same idea is:

\begin{verbatim}
value = 42

if 0 <= value and value <= 100:
    print("value is inside the accepted interval")
\end{verbatim}

The chained version is closer to the intended reading: the value is between the lower and upper limits.

For C programmers, this is a major semantic difference. In C, an expression that visually resembles this:

\begin{verbatim}
/* C idea, not Python */
0 <= value <= 100
\end{verbatim}

does not mean the same thing. C evaluates it roughly as:

\begin{verbatim}
/* C idea, not Python */
(0 <= value) <= 100
\end{verbatim}

The first comparison produces an integer-like truth result, usually \texttt{0} or \texttt{1}, and then that result is compared with \texttt{100}. That is not an interval test. Python's comparison chaining is designed to mean the interval test directly.

Python also supports mixed comparison operators in one chain:

\begin{verbatim}
a = 5
b = 5
c = 10

check = a == b < c

print(f"check: {check}")
\end{verbatim}

Possible output:

\begin{verbatim}
check: True
\end{verbatim}

This means:

\begin{verbatim}
a == b and b < c
\end{verbatim}

It does not mean:

\begin{verbatim}
(a == b) < c
\end{verbatim}

The shared middle operand is part of both neighboring comparisons.

Membership and identity comparisons can also be chained, although readability should control whether this is a good idea.

\begin{verbatim}
x = 42
items = [10, 42, 90]

check = 10 < x in items

print(f"check: {check}")
\end{verbatim}

Possible output:

\begin{verbatim}
check: True
\end{verbatim}

This means:

\begin{verbatim}
10 < x and x in items
\end{verbatim}

The chain is legal, but it may surprise readers. In most teaching and production code, simple numeric chains are clearer than clever mixed chains.

A more readable form is:

\begin{verbatim}
x = 42
items = [10, 42, 90]

check = 10 < x and x in items

print(f"check: {check}")
\end{verbatim}

Possible output:

\begin{verbatim}
check: True
\end{verbatim}

The practical rule is:

\begin{quote}
Use comparison chaining for natural mathematical or interval-like relations. Avoid chains that force the reader to stop and decode the grammar.
\end{quote}

The control-flow connection is direct. A chained comparison is one condition expression, but internally it contains several comparison edges. Python evaluates them from left to right and stops as soon as the chain cannot succeed. The final result is a \texttt{bool}, and that boolean selects the branch route.